% !PRINTED 626
\weberpage{626}
Man kann also jede gebrochene Funktion von $\theta$ durch eine ganze Funktion von $\theta$ ersetzen und erhält so den Satz:

1. \gesperrt{Jede Funktion $\zeta$ des Körpers $\Omega$ läßt sich auf eine einzige Weise in die Form setzen:}
\weq{6}{\zeta = x_0 + x_1\theta + x_2\theta^2 + \cdots + x_{n-1}\theta^{n-1},}
\gesperrt{deren Koeffizienten $x_0, x_1, x_2, \ldots, x_{n-1}$ rationale Funktionen von $z$ sind.}

Daß dies überhaupt möglich ist, schließt man aus (4) und (5); daß es nur auf eine Weise möglich ist, folgt aus der Irreduzibilität von $f(\theta)$; denn hätte man zwei verschiedene Darstellungen einer und derselben Größe $\zeta$ in der Form (6), so würde ihre Differenz eine Gleichung für $\theta$ von niedrigerem als $n$tem Grade ergeben.

Wählt man $n$ Funktionen des Körpers $\Omega$ beliebig aus:
\weq{7}{
\begin{aligned}
\eta_1 &= x_0^{(1)} + x_1^{(1)}\theta + \cdots + x_{n-1}^{(1)}\theta^{n-1}, \\
\eta_2 &= x_0^{(2)} + x_2^{(2)}\theta + \cdots + x_{n-1}^{(2)}\theta^{n-1}, \\% [sic] printed x_2^{(2)}, not x_1^{(2)}
&\qquad\qquad\qquad\cdots\cdots\cdots\cdots \\
\eta_n &= x_0^{(n)} + x_1^{(n)}\theta + \cdots + x_{n-1}^{(n)}\theta^{n-1},
\end{aligned}
}
mit der einzigen Beschränkung, daß die Determinante
\[
\Sigma \pm x_0^{(1)} x_1^{(2)} \cdots x_{n-1}^{(n)}
\]
nicht identisch Null ist, so kann man durch Elimination der Potenzen von $\theta$ aus (6) und (7) jede Funktion $\zeta$ in $\Omega$ in die Form setzen:
\weq{8}{\zeta = y_1\eta_1 + y_2\eta_2 + \cdots + y_n\eta_n,}
deren Koeffizienten $y_i$ rationale Funktionen von $z$ sind.

Ein solches System von Funktionen $\eta_1, \eta_2, \ldots, \eta_n$ heißt eine \gesperrt{Basis des Körpers} $\Omega$.

Damit ein Funktionensystem $\eta_1, \eta_2, \ldots, \eta_n$ eine Basis von $\Omega$ bilde, ist notwendig und hinreichend, daß keine Gleichung von der Form
\[
y_1\eta_1 + y_2\eta_2 + \cdots + y_n\eta_n = 0
\]
bestehe, in der die Koeffizienten $y_1, y_2, \ldots, y_n$ nicht sämtlich identisch verschwinden.

So bilden die Funktionen
\[
1, \theta, \theta^2, \ldots, \theta^{n-1}
\]
eine Basis von $\Omega$.
